% from ega4-2-fr.tex
\oldpage[IV-2]{39}
\begin{proposition}[3.1.10]
\label{IV.2.3.1.10-fr}
Soient $X$, $Y$ des préschémas localement noethériens, $f:X\to Y$ un
morphisme entier. Alors, pour tout $\mathcal O_X$-Module quasi-cohérent
$\mathcal F$, on a
$f(\operatorname{Ass}(\mathcal F))=\operatorname{Ass}(f_*(\mathcal F))$.
\end{proposition}

La question étant locale sur $Y$ et le morphisme $f$ étant affine, on est
immédiatement ramené au cas où $Y$ est affine, autrement dit au

\begin{lemma}[3.1.10.1]
\label{IV.2.3.1.10.1-fr}
Soient $A$, $B$ deux anneaux noethériens, $\rho:A\to B$ un homomorphisme
d'anneaux faisant de $B$ une $A$-algèbre entière, $M$ un $B$-module. Alors les
idéaux premiers $\mathfrak p\in\operatorname{Ass}(M_{[\rho]})$ sont les images
réciproques par $\rho$ des idéaux premiers
$\mathfrak q\in\operatorname{Ass}(M)$.
\end{lemma}

En effet, si $\mathfrak q\in\operatorname{Ass}(M)$, $\mathfrak q$ est
l'annulateur dans $B$ d'un élément $x\in M$, donc
$\rho^{-1}(\mathfrak q)$ est l'annulateur dans $A$ de $x$. Inversement, soit
$\mathfrak p\in\operatorname{Ass}(M_{[\rho]})$, de sorte que $\mathfrak p$
est l'image réciproque par $\rho$ de l'annulateur $\mathfrak b$ dans $B$ d'un
élément $x\in M$~; il résulte donc du premier théorème de Cohen-Seidenberg
qu'il existe un idéal premier $\mathfrak q$ de $B$ contenant $\mathfrak b$ et
dont l'image réciproque est $\mathfrak p$ (Bourbaki, \emph{Alg. comm.},
chap.~V, \S~2, n\textsuperscript{o}~1, cor.~2 du th.~1)~; en considérant un
des idéaux premiers minimaux parmi ceux contenus dans $\mathfrak q$ et
contenant $\mathfrak b$, on peut évidemment supposer que $\mathfrak q$
lui-même est un de ces idéaux minimaux. Mais comme $B\cdot x\subset M$ est
isomorphe à $B/\mathfrak b$, on sait que l'on a alors
$\mathfrak q\in\operatorname{Ass}(B/\mathfrak b)\subset\operatorname{Ass}(M)$
(Bourbaki, \emph{Alg. comm.}, chap.~IV, \S~1, n\textsuperscript{o}~4,
th.~2).

\begin{corollary}[3.1.11]
\label{IV.2.3.1.11-fr}
Sous les hypothèses de (3.1.10), pour que $\mathcal F$ soit sans cycle premier
associé immergé, il suffit qu'il en soit de même de $f_*(\mathcal F)$.
\end{corollary}

Supposons en effet que $f_*(\mathcal F)$ n'ait pas de cycle premier associé
immergé. Notons que si $A$ est une algèbre entière sur un corps $k$, tous les
idéaux premiers de $A$ sont maximaux (Bourbaki, \emph{Alg. comm.}, chap.~V,
\S~2, n\textsuperscript{o}~1, prop.~1)~; il résulte donc de (I, 6.2.2) que les
fibres de $f$ sont des espaces \emph{discrets}~; si $x$, $x'$ sont deux points
distincts de $\operatorname{Ass}(\mathcal F)$, aucun d'eux ne peut donc être
adhérent à l'autre si $f(x)=f(x')$~; et si $f(x)\ne f(x')$, (3.1.10) et
l'hypothèse entraînent qu'aucun des deux points $f(x)$, $f(x')$ ne peut être
adhérent à l'autre, donc il en est de même de $x$ et $x'$.

\begin{remark}[3.1.12]
\label{IV.2.3.1.12-fr}
Sous les hypothèses de (3.1.10), il peut se faire par contre que $\mathcal F$
soit sans cycle premier associé immergé, mais non $f_*(\mathcal F)$. Prenons
par exemple $Y=\operatorname{Spec}(k[T])$ où $k$ est un corps («~droite
affine~»), et $X$ somme de $X_1=Y$ et de $X_2=\operatorname{Spec}(k)$, le
morphisme $X_2\to Y$ correspondant à l'homomorphisme canonique
$k[T]\to k[T]/\mathfrak m$, où $\mathfrak m$ est l'idéal maximal $(T)$. Il
est clair que le morphisme $f:X\to Y$ est fini~; si l'on prend
$\mathcal F=\mathcal O_X$, $\mathcal F$ est sans cycle premier associé
immergé, mais $f_*(\mathcal F)=\widetilde M$, où $M$ est le $k[T]$-module
somme directe de $k[T]$ et de $k$, donc $\operatorname{Ass}(M)$ est formé du
point générique $(0)$ de $Y$ et du point $\mathfrak m$.
\end{remark}

\begin{proposition}[3.1.13]
\label{IV.2.3.1.13-fr}
Soient $X$ un préschéma localement noethérien, $U$ un ouvert de $X$,
$i:U\to X$ l'injection canonique. Pour tout $\mathcal O_U$-Module
quasi-cohérent $\mathcal F$, on a
$\operatorname{Ass}(i_*(\mathcal F))=\operatorname{Ass}(\mathcal F)$.
\end{proposition}

Rappelons que $i_*(\mathcal F)$ est un $\mathcal O_X$-Module quasi-cohérent
(I, 2.2 et 1.7.4)~; comme $i_*(\mathcal F)|U=\mathcal F$, on a
$(\operatorname{Ass}(i_*(\mathcal F)))\cap U=\operatorname{Ass}(\mathcal F)$,
et il reste donc à prouver que $\operatorname{Ass}(i_*(\mathcal F))\subset U$.
Mais d'après (3.1.8), cette relation signifie que pour tout ouvert affine $V$
de $X$, toute section de $i_*(\mathcal F)$ au-dessus de $V$ qui est nulle
dans $U\cap V$, est nulle, condition trivialement vérifiée puisque
$\Gamma(V,i_*(\mathcal F))=\Gamma(U\cap V,\mathcal F)
=\Gamma(U\cap V,i_*(\mathcal F))$.

